\documentclass[aspectratio=169,11pt]{beamer}
\usetheme{Madrid}
\usecolortheme{default}
\setbeamertemplate{navigation symbols}{}
\setbeamertemplate{footline}[frame number]

\usepackage[utf8]{inputenc}
\usepackage[T1]{fontenc}
\usepackage[italian]{babel}
\usepackage{lmodern}
\usepackage{amsmath,amssymb,mathtools}
\usepackage{bm}
\usepackage{array}


\usepackage{tikz}
\usetikzlibrary{calc}
\usetikzlibrary{positioning}


\newcommand{\astr}[1]{\textsf{{(\boldmath{${*}$}#1)}}}
\newcommand{\tr}{\operatorname{tr}}
\newcommand{\C}{\mathbb{C}}
\newcommand{\R}{\mathbb{R}}
\newcommand{\D}{\mathbb{D}}
\newcommand{\M}{\mathbb{M}^{2,1}}
\newcommand{\HH}{\mathbb{H}}
\newcommand{\QDD}{\mathbf{Q}_{\D}}
\newcommand{\EE}{\mathbf{E}}
\newcommand{\Amap}{\mathcal{A}}
\newcommand{\I}{\mathbf{I}}
\newcommand{\II}{\mathbf{II}}
\newcommand{\Iff}{\mathbf{I}}
\newcommand{\IIff}{\mathbf{II}}
\newcommand{\Hopf}{\operatorname{Hopf}}
\newcommand{\norm}[1]{\left\lVert #1 \right\rVert}


\newcommand{\placeholder}[1]{%
  \begin{center}
  \fbox{\parbox{0.82\textwidth}{\centering \vspace{1.2cm}#1\\[0.4em]\footnotesize (spazio riservato per un eventuale diagramma/figura)\vspace{1.2cm}}}
  \end{center}}




\title{Classificazione delle superfici space-like iperboliche a curvatura media costante in spazio di Minkowski}
\subtitle{Tesi di Laurea in Matematica}
\author{\centering \textbf{Candidato:} Elia Andrea Perli\\
		\textbf{Relatore:} Prof. Francesco Bonsante}
\institute{Università degli Studi di Pavia}
\date{  2026 - 04 - 28}










\begin{document}

\begin{frame}
  \titlepage
\end{frame}

\begin{frame}{Struttura della presentazione}
\tableofcontents[hideallsubsections]
\end{frame}

\section{Setup}


\begin{frame}{Setting geometrico}
	\begin{itemize}
		\item Lo spazio di Minkowski è $\M=\R^3$ con pseudometrica
			\[
				g_{\M} = ds^2 = (dx^1)^2+(dx^2)^2-(dx^3)^2.
			\]
		\item Una superficie immersa $ \iota:S\hookrightarrow \M$ è
			\alert{space-like} se la metrica indotta è Riemanniana.
			Questa è la prima forma fondamentale della superficie
			$\I = \iota^*(g_{\M})$. Si definisce analogamente al caso euclideo la seconda
			forma fondamentale $\II$.
		\item Se la superficie è space-like, questa è localmente della forma
			$x^3=f(x^1,x^2)$ con $|\nabla f|<1$.
	\end{itemize}
	\vspace{0mm}
	\begin{columns}[T,onlytextwidth]
		\begin{column}{0.48\textwidth}
			\begin{itemize}
				\item Valgono le equazioni di Gauss-Codazzi
					\begin{align*}
						\begin{cases}
							\det_{\I}(\II)=-K,\\
							d^{\nabla_\I}\II=0\ .
						\end{cases}
					\end{align*}
					Queste equazioni di compatibilità permettono di
					ricostruire la superficie.
			\end{itemize}
		\end{column}

		\begin{column}{0.48\textwidth}
			\begin{itemize}
				\item Come nel caso euclideo, da $\I$ e $\II$ si definisce
					la curvatura media:
					\begin{align*}
						H=\frac{1}{2}\tr_{\I}(\II).
					\end{align*}
			\end{itemize}
		\end{column}
	\end{columns}
\end{frame}

\begin{frame}{Superfici complete a curvatura media costante}
\begin{itemize}
	\setlength\itemsep{2mm}
	\item Consideriamo delle superfici space-like complete a curvatura
	media costante (CMC):
	\begin{align*}
		H=\text{cost}.
	\end{align*}
	\item Per una CMC, questa è un grafico intero $f(x^1,x^2) =x^3$ se e soltanto se la prima forma fondamentale è completa. In particolare è semplicemente connessa.
	\item Usando un risultato di Cheng e Yau $f$ è una funzione convessa. Inoltre se $H=0$, $S$ è un iperpiano.
	\item Se $H>0$, riscalando l'immersione possiamo normalizzare $H=1$.
\end{itemize}

\begin{block}{Obiettivo}
	Classificare le superfici space-like complete con $H=1$, a meno di
	isometria di $\M$.
\end{block}
\end{frame}


\begin{frame}{Parametrizzazione conforme}
\begin{itemize}
  \item Essendo $S$ una superficie riemanniana con $\I$ possiamo trovare localmente delle coordinate isoterme $(x,y)$, ovvero tali che
  \[
    \I=e^{2u}(dx^2+dy^2)=e^{2u}|dz|^2,
    \qquad z=x+iy.
  \]
  \item Il cambio di coordinate $z=x+iy$ determina una struttura complessa sulla superficie.
	\item $S$ è semplicemente connessa. Il teorema di uniformizzazione implica l'esistenza di un diffeomorfismo conforme
	\[
	\sigma:\Sigma \to S,
	\qquad
	\Sigma \in \{\C,\widehat{\C},\mathbb{D}\}.
	\]
  \item $S$ è iperbolica, dunque otteniamo una parametrizzazione
  \[
    \sigma:\D\to S.
  \]
  Dove $\D$ dotato della metrica $ ds_p^2=\frac{4}{(1-|z|^2)^2}|dz|^2$.
\end{itemize}

\section{Il risultato di classificazione, la mappa $\Amap$}
\end{frame}
\begin{frame}{Le incognite in coordinate isoterme}
\begin{itemize}
	\setlength\itemsep{2mm}

	\item Fissate coordinate isoterme su $\D$, scriviamo
	$\I=e^{2u}(dx^2+dy^2)$. Poiché $H=1$, la seconda forma fondamentale
	è della forma
	\[
		\II=e^{2u}(dx^2+dy^2)
		+h_{11}(dx^2-dy^2)+2h_{21}\,dxdy.
	\]

	\item Definiamo
	\[
		\varphi:=h_{11}-ih_{21}.
	\]
	L'equazione di Codazzi per $\II$ equivale al fatto che $\varphi$ sia
	olomorfa su $\D$.

	\item L'equazione di Gauss è equivalente alla PDE
	\begin{align}
		\Delta u=e^{2u}-|\varphi|^2e^{-2u}.
		\tag{\boldmath{$*$}} \label{eq:gausspde}
	\end{align}
\end{itemize}

\begin{center}
\begin{minipage}{0.72\textwidth}
\begin{block}{Riassumendo} 
	\begin{center}
		$(\I,\II)\longmapsto(\varphi,u)$ dove

		\vspace{1mm}
		\begin{minipage}{0.72\linewidth}
		\begin{itemize}
			\item $\varphi=h_{11}-ih_{21}$ è olomorfa;
			\item $u$ verifica l'equazione \eqref{eq:gausspde}.
		\end{itemize}
		\end{minipage}
	\end{center}
\end{block}
\end{minipage}
\end{center}
\end{frame}
\begin{frame}{Codazzi produce il dato olomorfo}
\begin{itemize}
	\setlength\itemsep{2mm}

	\item Viceversa, fissiamo una funzione olomorfa
	$\varphi:\D\to\C$ e definiamo
	\[
		h_{11}=\Re(\varphi),
		\qquad
		h_{21}=-\Im(\varphi).
	\]

	\item Con questa scelta l'equazione di Codazzi è automaticamente
	verificata.

	\item Rimane da scegliere $u$ in modo che l'equazione di Gauss sia
	soddisfatta e che la metrica
	\[
		\I=e^{2u}(dx^2+dy^2)
	\]
	sia completa.
\end{itemize}

Abbiamo ridotto la classificazione seguente problema.
\begin{center}
\vspace{-2mm}
\begin{minipage}{0.72\textwidth}
\begin{block}{Problema}
		Data $\varphi$ olomorfa, trovare $u$ tale che
		\[
			\Delta u=e^{2u}-|\varphi|^2e^{-2u}
		\]
		e
		$\I=e^{2u}(dx^2+dy^2)$ sia completa.
\end{block}
\end{minipage}
\end{center}
\end{frame}
\begin{frame}{Riassumendo}
\begin{itemize}
	\item Abbiamo definito una relazione che associa ad una CMC una funzione olomorfa sul disco $\mathbb{D}$, \[ S \mapsto \varphi .\] 
	\item Quozientando per le relazioni di isometria $S \sim S'$ e trasformazioni conformi di $\mathbb{D}$, $\varphi \sim \varphi'$, otteniamo la mappa $\Amap$
		\[\Amap : \mathbf{E} \to \QDD\ ,\] 

	Dove $\EE$ è lo spazio delle superfici CMC-$1$ iperboliche a meno di isometria, e $\QDD$ lo spazio dei differenziali quadratici olomorfi su $\D$, a meno di cambio di trasformazioni conformi di $\mathbb{D}$.
\end{itemize}
\vspace{2mm}
\begin{center}
\begin{minipage}{0.68\textwidth}
\begin{block}{Teorema di classificazione}
	\begin{center}
		L'applicazione $\Amap$ è ben definita e biettiva.
	\end{center}
\end{block}
\end{minipage}
\end{center}

\vspace{1mm}
\end{frame}





\begin{frame}[t]{Il problema analitico}
	Abbiamo ridotto il teorema di classificazione alla risoluzione del seguente problema che assume la seguente forma perché l'equazione può essere riformulata nella nuova incognita $w$ tale che $e^{2u}|dz|^2 = e^{2w} ds_p^2$.

  \begin{columns}[T,onlytextwidth]
    \column{0.08\textwidth}\hfill
    \column{0.84\textwidth}
    \begin{block}{ Problema }
      Dato $\Phi\in\QDD$, esiste \alert{un'unica}
      $w\in C^\infty(\D)$ che verifica
      \[
        \Delta_p w \;=\; e^{2w}-\norm{\Phi}^2 e^{-2w}-1
        \quad\text{su } \D,
      \]
	  e tale che
      \begin{itemize}\itemsep=2pt
        \item $e^{2w}\,ds_p^2$ è una metrica completa su $\D$;
        \item $e^{2w}-\norm{\Phi}^2 e^{-2w} > 0$ su $\D$.    
      \end{itemize}
    \end{block}
    \column{0.08\textwidth}\hfill
  \end{columns}

  \vspace{3mm}
  \begin{itemize}
    \item \alert{esistenza}: metodo delle sotto-soluzioni/sovra-soluzioni.
    \item \alert{unicità}: principio del massimo generalizzato;
  \end{itemize}
\end{frame}

\begin{frame}[t]{Esistenza: sotto/sovra-soluzioni ed esaustione}

\begin{block}{Metodo delle sotto/sovra-soluzioni}
Sia \(M\) una varietà Riemanniana completa e sia \(F(x,u)\in C^\infty(M\times\R)\) con \(\partial_uF>0\).
Se esistono \(\underline\psi,\overline\psi\in C^0(M)\cap H^1(M)\) con \(\underline\psi\le \overline\psi\) tali che, debolmente,
\(\Delta \overline\psi \le F(x,\overline\psi)\) e \(\Delta \underline\psi \ge F(x,\underline\psi)\),
allora esiste \(u\in C^\infty(M)\) tale che \(\Delta u=F(x,u)\) e \(\underline\psi\le u\le \overline\psi\).
\end{block}

\vspace{1mm}

Nel nostro caso:
\begin{itemize}
	\item Si risolve il problema nel caso limitato esibendo una sottosoluzione e una sovrasoluzione.
  \item Nel caso generale, si costruisce una successioni di limitati $(D_k)_k$ che invade \(\D\). Su ciascun dominio $\|\Phi\|$ è limitata.  Da questo si trova poi la convergeza in $C^2_{\mathrm{loc}}$
\end{itemize}

\end{frame}



\section{Il problema analitico}



\begin{frame}[t]{Unicità: principio del massimo generalizzato}
  \begin{columns}[T,onlytextwidth]
    \column{0.07\textwidth}\hfill
    \column{0.86\textwidth}
    \begin{block}{Principio del massimo generalizzato (Omori-Yau) }
      Sia $(M,g)$ completa, semplicemente connessa, con curvatura sezionale
      non positiva e curvatura di Ricci limitata dal basso. Se
      $f\in C^2(M)$ è limitata superiormente, allora esiste una successione
      $(p_n)\subset M$ tale che
      \[
        f(p_n)\to \sup_M f,\qquad
        |\nabla_g f|(p_n)\to 0,\qquad
        \limsup_n \Delta_g f(p_n)\le 0.
      \]
    \end{block}
    \column{0.07\textwidth}\hfill
  \end{columns}

  \vspace{2mm}
  Nella dimostrazione di unicità:
  \begin{itemize}
    \item si fissano due soluzioni $w,\tilde w$;
    \item si considera la loro differenza $\eta=w-\tilde w$;
    \item si mostra che $\eta$ è limitata superiormente;
    \item si applica il principio del massimo generalizzato a $\eta$ su $(\mathbb{D},\I)$.
  \end{itemize}
\end{frame}
\begin{frame}[t]{Unicità: conclusione}

\begin{itemize}
  \item Otteniamo una successione $(z_k)_k\subset\D$ tale che
  \[
    \eta(z_k)\to \bar\eta:=\sup_\D\eta,
    \qquad
    \limsup_{k\to\infty}\Delta_\I\eta(z_k)\leq 0.
  \]

  \item L'equazione soddisfatta da $w$ e $\tilde w$ implica, dove $|\mu|^2 = \|\Phi\|^2 e^{-4w}$
  \[
    \Delta_\I\eta
    =
    e^{2\eta}-|\mu|^2e^{-2\eta}-1+|\mu|^2.
  \]

  \item A meno di sottosuccessioni possiamo
  assumere
  \[
    |\mu|^2(z_k)\to a,
    \qquad
    0\leq a\leq 1.
  \]
  Passando al limite lungo $(z_k)_k$ otteniamo
  \[
    0
    \geq
    e^{2\bar\eta}-ae^{-2\bar\eta}-1+a.
  \]

  \item Definendo
  \[
    \psi(s)=e^{2s}-ae^{-2s}-1+a,
  \]
  si ha $\psi'(s)>0$ e $\psi(0)=0$. Dunque
  $\psi(\bar\eta)\leq 0$ implica $\bar\eta\leq 0$.

  \item Quindi $w\leq \tilde w$. Scambiando i ruoli delle due soluzioni,
  otteniamo $w=\tilde w$.
\end{itemize}
\end{frame}
\begin{frame}
  \vfill
  \begin{center}
    \Huge Grazie per l'attenzione.
  \end{center}
  \vfill
\end{frame}
\begin{frame}[allowframebreaks]{Riferimenti essenziali}
\footnotesize
\begin{thebibliography}{9}
\bibitem{Wan} T. Y. H. Wan, \emph{Constant mean curvature surface, harmonic maps, and universal Teichmüller space}, J. Differential Geom. 35 (1992), 643--657.
\bibitem{Jost} J. Jost, \emph{Compact Riemann Surfaces}, Springer, 2006.
\bibitem{doCarmo} M. P. do Carmo, \emph{Differential Geometry of Curves and Surfaces}, Dover, 2016.
\end{thebibliography}
\end{frame}

\end{document}
